\section{Logique de sélection des projets}

Le dossier s'appuie sur une sélection de projets professionnels, académiques et complémentaires. Le projet \terme{fcu} / Ediacara constitue le projet principal, car il permet de démontrer une démarche complète : analyse d'un existant complexe, adaptation technique, tests, documentation et transmission de connaissances.

Les autres projets complètent la couverture du référentiel, notamment sur le développement back-end, le développement front-end, le traitement de données, le \terme{devops} et le pilotage d'équipe. Le homelab personnel occupe une place volontairement complémentaire : il ne constitue pas une mission professionnelle ni un projet académique évalué, mais un environnement d'expérimentation continue utilisé pour approfondir les pratiques de déploiement, de GitOps et d'observabilité abordées en formation.

\section{Projets retenus}

Les projets retenus sont les suivants : \terme{fcu} / Ediacara, \terme{s123}, FeastVerse, Brise-glace, Projet Data / IA, Homelab personnel : GitOps, CI/CD et observabilité, et IPF 2026 Lynx Immo.

\section{Correspondance générale avec les blocs}

Une synthèse détaillée par compétence est présentée en fin de dossier. Cette cartographie initiale sert de repère de lecture pour comprendre le rôle de chaque projet dans la validation des blocs de compétences.
